\noindent\textbf{Purpose:} Training on facial-expression vote distributions may improve distributional agreement without improving confidence-selected category predictions. We examine how that distinction depends on target probability values, tied maxima, accepted images and output vocabulary.
\par\smallskip
\noindent\textbf{Methods:} An adaptive within-source study on FER2013/FER+ comprised 24 original Swin-T/ResNet-18 runs and twelve additional Swin-T runs, with three paired seeds per condition. Soft retained vote proportions; Uniform kept one dominant probability and flattened the remainder, whereas Tie-Uniform preserved all tied maxima. A permutation preserved each target's probability multiset. Seven- and eight-expression classifiers were scored against all ten annotation categories. The primary reporting point accepted each model's 80\% most confident images. Analyses included fixed-image comparisons, checkpoint rules and confidence scores. Frozen original models were evaluated on 1,440 RAVDESS clips.
\par\smallskip
\noindent\textbf{Results:} Distance-selected seven-output Soft reduced disagreement by 0.79 percentage points relative to the entropy-preserving permutation. With eight outputs, Soft achieved 23.74\% disagreement versus 24.58\% for Uniform and 23.79\% for Tie-Uniform. The Soft-minus-Tie-Uniform difference was $-0.05$ points, with a conditional 95\% interval of $[-0.31,+0.21]$, and reversed under alternative rankings. Soft had lower squared distance in every matched softening-control comparison, but the original selective advantage did not persist across backbones. External transfer was poor and sensitive to aggregation and annotation reference.
\par\smallskip
\noindent\textbf{Conclusion:} The measured benefit of crowd-distribution supervision depends on the control and selection policy. Better distributional agreement does not establish a general advantage in selective categorical prediction.
